\documentclass[11pt,a4paper]{article}
\usepackage[T1]{fontenc}
\usepackage{isabelle,isabellesym}
\usepackage{amsmath,amssymb}
\usepackage{pdfsetup}
\isabellestyle{it}

\begin{document}
\title{Bounds for Finite Families of Complex Embeddings}
\author{OpenAI Codex}
\maketitle

\begin{abstract}
We develop maximum-modulus bounds for finite families of complex-valued
maps, recover basis coordinates through an indexed inverse basis matrix,
and lift bounded integer kernels to complex kernels using integral
structure constants. The lifted results include Hadamard, linear, and
Siegel-exponent bounds. A two-point family models the conjugate images of
the Gaussian basis $1,i$. The session depends on the AFP entry Bounded
Integer Kernels and abstracts techniques in Mathlib's number-field house
and Siegel-lemma developments~\cite{mathlib-house,mathlib-siegel}.
\end{abstract}

\section{Scope and main bounds}
The finite set $E$ labels complex-valued maps. The general results use
pointwise complex values and do not require a particular number field.
The \emph{house} of a family $z$ is the maximum of $|z(e)|$ for $e\in E$,
with value zero when $E$ is empty. An indexed biorthogonal family provides
coordinate recovery and an inverse basis matrix.

Suppose a $p\times q$ complex matrix acts on a rank-$d$ integer span,
with $0<p<q$ and $d>0$. If its integral structure constants have absolute
value at most $M$, the lifted Siegel lemma gives a nonzero kernel vector
whose integer coordinates have absolute value at most
\[
  \bigl(qd\,\max(1,M)\bigr)^{pd/(qd-pd)}.
\]
If each basis element has house at most $K$, every coordinate of the
resulting complex kernel vector, across the finite family, has house at
most $dK$ times this bound. The theories also retain determinant and
linear bounds for applications needing those particular estimates.
The three estimates share one witness property: the vector is nonzero and
lies in the kernel, it has integral basis coordinates bounded by the stated
constant, and both its house and its modulus are bounded by $dK$ times
that constant. A separate named bound estimates structure constants by
$\lceil |E|RHK\rceil$ from recovery coefficients of house at most $R$
and matrix entries of house at most $H$.
The lifting theorems are organized in two locales. The structure-constant
locale collects the matrix, basis, multiplication table, and house hypotheses.
The entry locale adds coordinate recovery and entry-house bounds and
interprets the structure-constant locale. The determinant, exponent, and
linear results then differ only in their rank conditions and quantitative
bounds.

\section{A two-point example}
The final theory uses two labels for the identity and conjugate images
of $1$ and $i$. It constructs the dual basis explicitly, verifies
biorthogonality, recovers coordinates, and proves integer independence
at the identity image. This illustrates the hypotheses of the abstract
recovery interface without invoking a general number-field library.

\input{session}

\begin{thebibliography}{9}
\bibitem{mathlib-house}
Michail Karatarakis, \emph{House of an Algebraic Number},
Mathlib, \texttt{Mathlib/NumberTheory/NumberField/House.lean}, 2024.

\bibitem{mathlib-siegel}
Fabrizio Barroero, Laura Capuano, and Amos Turchet,
\emph{Siegel's Lemma}, Mathlib,
\texttt{Mathlib/NumberTheory/SiegelsLemma.lean}, 2024.
\end{thebibliography}
\end{document}
